%% air_freight_routing.tex
%% Air Freight Multi-Commodity Routing — MVP Formal Model
%% Phase 1 — must be approved before any code starts
%% Status: Draft v1 — 2026-05-16

\documentclass[11pt,a4paper]{article}

\usepackage{amsmath,amssymb,amsthm}
\usepackage{booktabs}
\usepackage{geometry}
\usepackage{hyperref}
\usepackage{array}
\usepackage{xcolor}
\usepackage{enumitem}
\usepackage{longtable}
\usepackage{graphicx}

\geometry{margin=2.5cm}
\setlength{\parskip}{0.4em}
\hypersetup{colorlinks=true, linkcolor=blue}

\newcommand{\ceil}[1]{\left\lceil #1 \right\rceil}
\newcommand{\floor}[1]{\left\lfloor #1 \right\rfloor}
\newcommand{\defeq}{\triangleq}

\theoremstyle{definition}
\newtheorem{defn}{Definition}
\newtheorem{remark}{Remark}

\title{\textbf{Air Freight Multi-Commodity Routing}\\[0.4em]
\large MVP Formal Model --- Phase 1, v1}
\author{AI Freight Agent}
\date{Draft --- 2026-05-16}

\begin{document}
\maketitle

\begin{abstract}
Formal specification of the optimization model for the air freight multi-commodity routing
problem. The model is a Binary Multi-Commodity Flow problem on a commodity-specific subgraph
of the air freight network (see \texttt{docs/air\_freight\_multi\_shipment\_graph.drawio}
and \texttt{docs/air\_freight\_shipment\_subgraph.drawio}). Scope: single-shipment routing
across multiple flight legs (direct and multi-hop) with two air capacity layers — contracted
ULD allocation (Block Space Agreements) and spot rate-card pricing. Probabilistic transit
time objectives, through-ULD as a decision variable, period-level utilization commitments,
and full bin-packing inside ULDs are deferred. The model parallels the Ocean FCL formulation
(\texttt{model/ocean\_fcl\_routing.tex}) for notational consistency.
\end{abstract}

\tableofcontents
\newpage

%%--------------------------------------------------------------------
\section{Problem Statement}
%%--------------------------------------------------------------------

Given a set of air shipment requests (commodities), each defined by origin location, destination
location, cargo volume, chargeable weight, and time constraints, find a minimum-cost assignment
of commodities to paths through the multimodal air freight network such that:
\begin{itemize}[noitemsep]
  \item every commodity is delivered on or before its deadline,
  \item no flight leg exceeds its ULD physical capacity (volume and weight per ULD),
  \item no Block Space Agreement (BSA) exceeds its contracted per-flight ULD allocation,
  \item period-level allocation caps are respected,
  \item cargo cutoff times are honored,
  \item hub minimum connect times are honored, including the through-ULD vs.\ re-ULD distinction, and
  \item optional per-commodity budget caps are respected.
\end{itemize}

\paragraph{MVP scope.}
\begin{itemize}[noitemsep]
  \item Single-shipment routing; each shipment routes independently subject to shared supply.
  \item Each scheduled flight leg is one arc. Multi-stop flights (same flight number, intermediate
        stops) are decomposed into multiple arcs for transit time visibility.
  \item ULD types: LD3, LD7, PMC pallet, AKE. Optimizer selects ULD type per shipment per flight.
  \item Two air supply layers per flight: (a) contracted ULD allocation under one or more BSAs;
        (b) spot rate-card at IATA weight breaks.
  \item Pivot weight (take-or-pay minimum) modeled on hard BSA contracts.
  \item Through-ULD vs.\ re-ULD at hub modeled via pre-computed compatibility parameter
        $\psi$, with rehandling cost and MCT differentiated. Explicit decision variable for
        through-ULD selection deferred to P1.
  \item Period-level minimum utilization commitments (hard BSA take-or-pay across period)
        are computed by a separate upstream model and surfaced as fixed per-run targets.
        This MILP does not solve them jointly.
  \item Objective: minimize total air + ground freight cost subject to deterministic
        time-window feasibility. Probabilistic on-time delivery objective deferred to P1.
  \item Hazmat / perishable / valuable cargo type segregation: model parameter, not a hard MVP
        constraint (recorded for future restriction logic).
\end{itemize}

%%--------------------------------------------------------------------
\section{Graph Structure}
%%--------------------------------------------------------------------

\subsection{Physical Network Graph \texorpdfstring{$G(N, A)$}{G(N,A)}}

The air freight network is a directed graph $G = (N, A)$.

\[
  N = N_O \cup N_{\text{CFS,O}} \cup N_{\text{AP,O}} \cup N_{\text{AP,H}} \cup N_{\text{AP,D}} \cup N_{\text{CFS,D}} \cup N_D
\]

\begin{description}[leftmargin=3.4cm, labelwidth=3.2cm, style=nextline]
  \item[$N_O$ — Origins] One per shipper origin door (e.g., factory, warehouse). Commodity source nodes.
  \item[$N_{\text{CFS,O}}$ — Origin CFS] Container Freight Stations near origin airports
        (e.g., TPE-CFS, KHH-CFS). ULD build-up location. Airline-owned empty ULDs are
        delivered here; cargo from one or more shipments is stuffed into a ULD; built ULD
        is trucked to the airline cargo terminal.
  \item[$N_{\text{AP,O}}$ — Origin airports] Airline cargo terminals at departure airports
        (e.g., TPE, KHH). Built ULDs are received here for loading onto aircraft.
  \item[$N_{\text{AP,H}}$ — Hub airports] Intermediate airports on multi-hop routes
        (e.g., HKG, NRT, ANC). Through-cargo handling or re-ULDing occurs here.
  \item[$N_{\text{AP,D}}$ — Destination airports] Airline cargo terminals at arrival airports
        (e.g., JFK, ORD, LAX). ULDs are unloaded here.
  \item[$N_{\text{CFS,D}}$ — Destination CFS] CFS near destination airport for ULD breakdown.
        Trucked from airline cargo terminal; ULDs broken down; individual consignments
        released for final delivery.
  \item[$N_D$ — Destinations] One per consignee delivery door. Commodity sink nodes.
\end{description}

\[
  A = A_{\text{PU}} \cup A_{\text{OC}} \cup A_{\text{AIR}} \cup A_{\text{DC}} \cup A_{\text{FD}}
\]

\begin{description}[leftmargin=3.4cm, labelwidth=3.2cm, style=nextline]
  \item[$A_{\text{PU}}$ — Pickup] $(i,j)$, $i \in N_O$, $j \in N_{\text{CFS,O}}$.
        Truck pickup from origin door to origin CFS.
  \item[$A_{\text{OC}}$ — Origin CFS to origin airport] $(i,j)$, $i \in N_{\text{CFS,O}}$,
        $j \in N_{\text{AP,O}}$. Built ULDs trucked from CFS to airline cargo terminal.
        Includes fixed CFS build-up time as part of effective arc duration.
  \item[$A_{\text{AIR}}$ — Flight legs] $(i,j)$, $i, j \in N_{\text{AP,O}} \cup N_{\text{AP,H}}
        \cup N_{\text{AP,D}}$. One arc per scheduled flight leg (specific carrier, flight
        number, departure date, origin airport, destination airport).
        A multi-stop flight (e.g., $f$: TPE$\to$ANC$\to$JFK) decomposes into two air arcs
        sharing the same flight number metadata.
  \item[$A_{\text{DC}}$ — Destination airport to destination CFS] $(i,j)$, $i \in N_{\text{AP,D}}$,
        $j \in N_{\text{CFS,D}}$. Trucked from airline cargo terminal to CFS.
        Includes fixed CFS breakdown time.
  \item[$A_{\text{FD}}$ — Final delivery] $(i,j)$, $i \in N_{\text{CFS,D}}$, $j \in N_D$.
        Trucked from destination CFS to consignee door.
\end{description}

\begin{remark}[Origin and destination CFS as explicit nodes]
Unlike ocean (where the POL terminal is one node), air introduces a mandatory CFS step at
origin (ULD build-up) and destination (ULD breakdown). Modeling CFS as an explicit node
permits: (a) the fixed CFS dwell time to enter feasibility checks, (b) future modeling of
CFS dock capacity, ULD availability at CFS, and (c) alternative-CFS routing (e.g., a
Kaohsiung shipment trucked to TPE-CFS for consolidation).
\end{remark}

\begin{remark}[Flight legs vs.\ flight numbers]
A real-world flight number (e.g., 5Y9123 TPE$\to$ANC$\to$JFK) operates one aircraft sortie
with an intermediate stop. We model this as \emph{two arcs} in $A_{\text{AIR}}$ — $(\text{TPE},
\text{ANC})$ and $(\text{ANC}, \text{JFK})$ — both tagged with the same flight number. This
preserves transit time visibility at each leg and provides the structural hook for per-leg
ULD allocation accounting. The fact that the ULD physically stays on the aircraft across the
intermediate stop is captured by the through-ULD compatibility parameter $\psi$
(Section~\ref{sec:through-uld}).
\end{remark}

\paragraph{Figure reference.}
See the rendered network diagrams: \texttt{docs/air\_freight\_multi\_shipment\_graph.drawio}
(multi-shipment view, supply sharing) and \texttt{docs/air\_freight\_shipment\_subgraph.drawio}
(single-shipment subgraph $G(N_k, A_k)$ for S3 TPE$\to$NYC). Render the .drawio sources to PDF
via drawio export before final approval.

\subsection{Commodity-Specific Subgraph \texorpdfstring{$G(N_k, A_k)$}{G(Nk,Ak)}}
\label{sec:subgraph}

The MILP for commodity $k$ operates on a subgraph $G(N_k, A_k) \subseteq G(N, A)$ containing
only arcs that (i) are time-feasible for commodity $k$, (ii) lie on at least one complete
path from $o(k)$ to $d(k)$, and (iii) are compatible with the commodity's cargo type and
ULD constraints (e.g., hazmat-restricted flights excluded for non-DGR-certified carriers).
See Section~\ref{sec:prefilter} for the construction algorithm.

%%--------------------------------------------------------------------
\section{Sets and Indices}
%%--------------------------------------------------------------------

\begin{center}
\begin{tabular}{lp{11.5cm}}
\toprule
Symbol & Description \\
\midrule
$K$ & Set of commodities (shipments), indexed by $k$ \\
$N,\, A$ & All nodes and arcs \\
$A_{\text{PU}},\, A_{\text{OC}},\, A_{\text{AIR}},\, A_{\text{DC}},\, A_{\text{FD}}$ & Arc subsets by type \\
$A^k \subseteq A$ & Feasible arc set for commodity $k$ (Section~\ref{sec:prefilter}) \\
$A^k_S \defeq A^k \cap A_S$ & Feasible arcs of type $S$ for commodity $k$ \\
$N_k$ & Nodes incident to $A^k$ \\
$F$ & Set of scheduled flight legs, indexed by $f$. Each flight is one arc in $A_{\text{AIR}}$. \\
$F^k \subseteq F$ & Flights corresponding to arcs in $A^k_{\text{AIR}}$ \\
$f(i,j)$ & Mapping: air arc $(i,j) \in A_{\text{AIR}}$ to its flight $f \in F$ \\
$U$ & Set of ULD types: $U = \{\text{LD3, LD7, PMC, AKE}\}$ \\
$U_f \subseteq U$ & ULD types available on flight $f$ (depends on aircraft type) \\
$C$ & Set of ULD allocation contracts (BSAs), indexed by $c$ \\
$C_f \subseteq C$ & Contracts that allocate ULDs on flight $f$ \\
$T$ & Set of allocation time periods (e.g., monthly), indexed by $t$ \\
$\tau_k(i)$ & Earliest arrival time of commodity $k$ at node $i$ \\
$\text{Hub}(j)$ & Set of consecutive air-arc pairs at hub node $j \in N_{\text{AP,H}}$ for shipment $k$: $\{((i,j), (j,\ell)) : (i,j), (j,\ell) \in A^k_{\text{AIR}}\}$ \\
\bottomrule
\end{tabular}
\end{center}

%%--------------------------------------------------------------------
\section{Parameters}
%%--------------------------------------------------------------------

\subsection{Commodity Parameters}

\begin{center}
\begin{tabular}{llp{6.5cm}l}
\toprule
Parameter & Symbol & Description & Unit \\
\midrule
Volume & $v_k$ & Cargo volume (sum of all pieces) & CBM \\
Actual weight & $w_k$ & Gross weight & kg \\
Chargeable weight & $cw_k$ & $\max(w_k,\, v_k \cdot 167)$ & kg \\
Piece dimensions & $D_k$ & Max single-piece L$\times$W$\times$H (for ULD fit check) & cm \\
Cargo ready & $t_k^{\text{rdy}}$ & Earliest pickup time from origin door & datetime \\
Deadline & $T_k^{\text{dead}}$ & Latest acceptable arrival at destination & datetime \\
Origin & $o(k)$ & Source node $\in N_O$ & — \\
Destination & $d(k)$ & Sink node $\in N_D$ & — \\
Cargo type & $\tau_k$ & General / DGR / PER / VAL / HEA & categorical \\
HS code & $\text{hs}_k$ & 6-digit Harmonized System code & — \\
Incoterm & $\text{inco}_k$ & EXW / FCA / CIF / DDP & — \\
Budget cap & $B_k$ & Hard cost ceiling; $\infty$ if not set & USD \\
\bottomrule
\end{tabular}
\end{center}

\paragraph{Chargeable weight.}
\begin{equation}
  cw_k \defeq \max\!\left(w_k,\ v_k \cdot 167\right)
  \label{eq:chargeable}
\end{equation}
IATA standard: $1\,\text{m}^3 = 167\,\text{kg}$ volumetric equivalent. Dense cargo (machinery,
electronics) rated on $w_k$; bulky-light cargo (furniture, packaging) rated on volumetric.

\subsection{Ground Arc Parameters (Pickup, CFS, Final Delivery)}

For each ground arc $(i,j) \in A_{\text{PU}} \cup A_{\text{OC}} \cup A_{\text{DC}} \cup A_{\text{FD}}$:

\begin{center}
\begin{tabular}{llp{6.5cm}l}
\toprule
Parameter & Symbol & Description & Unit \\
\midrule
Cost & $c_{ij}^{\text{G}}$ & Truck cost per shipment & USD \\
Mean transit & $\mu_{ij}^{\text{G}}$ & Mean transit time & days \\
Std dev & $\sigma_{ij}^{\text{G}}$ & Std dev of transit & days \\
CFS build-up time & $\beta^{\text{O}}$ & Fixed origin CFS dwell (ULD build-up). Configurable. & hours \\
CFS breakdown time & $\beta^{\text{D}}$ & Fixed destination CFS dwell (ULD breakdown). Configurable. & hours \\
\bottomrule
\end{tabular}
\end{center}

The effective transit time on the CFS$\to$airport arc $(i,j) \in A_{\text{OC}}$ includes the
CFS build-up: $\tilde{\mu}_{ij}^{\text{OC}} = \mu_{ij}^{\text{OC}} + \beta^{\text{O}}$. Likewise for
$A_{\text{DC}}$ with $\beta^{\text{D}}$.

\subsection{Flight (Air Arc) Parameters}
\label{sec:flight-params}

For each flight $f \in F$ (equivalently, each $(i,j) \in A_{\text{AIR}}$):

\begin{center}
\begin{tabular}{llp{7cm}l}
\toprule
Parameter & Symbol & Description & Unit \\
\midrule
Carrier & $\text{carrier}(f)$ & IATA carrier code (CI, BR, CX, JL, NH, 5Y, ...) & — \\
Flight number & $\text{fno}(f)$ & Flight number; multiple legs may share a number & — \\
Aircraft type & $\text{ac}(f)$ & 777F, 747-8F, A330F, 777-300ER, A350-900, ... & — \\
ETD & $\text{ETD}_f$ & Scheduled departure datetime & datetime \\
ETA & $\text{ETA}_f$ & Scheduled arrival datetime & datetime \\
Block time & $\mu_f^{\text{AIR}}$ & $\text{ETA}_f - \text{ETD}_f$ & hours \\
Std dev & $\sigma_f^{\text{AIR}}$ & Std dev of actual block time & hours \\
Cargo cutoff & $\text{CGC}_f$ & Latest time ULDs must be tendered to airline cargo terminal & datetime \\
DGR certified & $\text{dgr}(f)$ & Accepts dangerous goods? & boolean \\
Cold chain & $\text{per}(f)$ & Maintains $\leq 8^\circ$C? & boolean \\
\bottomrule
\end{tabular}
\end{center}

\subsection{ULD Specifications}
\label{sec:uld-spec}

\begin{center}
\begin{tabular}{lllll}
\toprule
ULD type & Code & Volume & Payload & Aircraft compatibility \\
\midrule
LD3       & AKE/AVE & 4.5 m\textsuperscript{3} & 1{,}587 kg & Wide-body belly + freighter \\
LD7       & ALE/AAY & 11.1 m\textsuperscript{3} & 4{,}626 kg & Wide-body belly + freighter \\
PMC pallet & PMC/PAG & 7.5 m\textsuperscript{3} & 6{,}804 kg & Freighter main deck \\
AKE small  & AKE     & 4.5 m\textsuperscript{3} & 1{,}497 kg & Wide-body belly \\
\bottomrule
\end{tabular}
\end{center}

For each $u \in U$:
\begin{center}
\begin{tabular}{llp{6.5cm}l}
\toprule
Parameter & Symbol & Description & Unit \\
\midrule
Volume capacity & $V_u$ & ULD usable volume (Table above) & CBM \\
Weight capacity & $W_u$ & ULD payload limit & kg \\
Fill rate cap & $\eta$ & Max fraction of CBM usable in practice (default $0.97$) & — \\
\bottomrule
\end{tabular}
\end{center}

For each flight $f$ and ULD type $u \in U_f$:
\begin{center}
\begin{tabular}{llp{6.5cm}l}
\toprule
Parameter & Symbol & Description & Unit \\
\midrule
Aircraft ULD positions & $P_{f,u}$ & Number of $u$-type ULD positions on $\text{ac}(f)$ & — \\
\bottomrule
\end{tabular}
\end{center}

$P_{f,u}$ is an upper bound on physical positions for ULD type $u$ on the aircraft. It is
independent of the forwarder's contracted allocation.

\subsection{Contracted ULD Allocation (Block Space Agreement) Parameters}
\label{sec:bsa-params}

A Block Space Agreement (BSA) reserves a fixed allocation of ULD positions per flight at a
contracted rate. For each contract $c \in C$:

\begin{center}
\begin{tabular}{llp{6.5cm}l}
\toprule
Parameter & Symbol & Description & Unit \\
\midrule
Flights covered & $F_c \subseteq F$ & Set of scheduled flights covered by contract $c$ & — \\
Per-flight allocation & $\text{alloc}(c, f, u)$ & ULD count of type $u$ allocated on flight $f$ under contract $c$ & ULDs \\
Period allocation & $\text{alloc}(c, u, t)$ & Total ULDs of type $u$ allocated in period $t$ under contract $c$ & ULDs \\
Period utilization & $\text{util}(c, u, t)$ & Already-booked ULDs in period $t$ at solve time & ULDs \\
Remaining (period) & $\text{rem}(c, u, t)$ & $\text{alloc}(c,u,t) - \text{util}(c,u,t)$ & ULDs \\
Min utilization (period) & $\underline{M}_{c,u,t}$ & Hard BSA minimum: external upstream model output. Default $0$ for soft BSA. & ULDs \\
Pivot weight & $\pi_{c,u}$ & Minimum chargeable weight per ULD per flight under $c$ & kg/ULD \\
Contract rate & $r_c$ & Contracted price & USD/kg chargeable \\
Hard vs.\ soft & $\text{hard}_c$ & Boolean: hard BSA (take-or-pay) or soft (cancellable) & — \\
\bottomrule
\end{tabular}
\end{center}

\begin{remark}[Period commitments are upstream]
$\underline{M}_{c,u,t}$ is computed by a separate upstream model that allocates how much of
each hard BSA's monthly take-or-pay commitment should be met by each routing run within the
period. The routing MILP takes $\underline{M}_{c,u,t}$ as a fixed input per run. The MILP
does \emph{not} simultaneously optimize period-level commitment allocation; that is a
higher-level decision.
\end{remark}

\begin{remark}[Pivot weight semantics]
The pivot $\pi_{c,u}$ is the minimum chargeable weight per ULD per flight regardless of actual
cargo loaded. If a forwarder claims one LD3 with pivot $1{,}000\,\text{kg}$ on a flight but
loads only $400\,\text{kg}$ of cargo, the forwarder pays for $1{,}000\,\text{kg} \cdot r_c$.
This is the take-or-pay floor on hard BSAs.
\end{remark}

\subsection{Spot Rate-Card Parameters}
\label{sec:spot-params}

For each flight $f$, the airline offers spot rate-card pricing with IATA weight breaks:

\begin{center}
\begin{tabular}{llp{6.5cm}l}
\toprule
Parameter & Symbol & Description & Unit \\
\midrule
Break tiers & $B = \{N, 45, 100, 300, 500, 1000\}$ & IATA standard break points (kg) & kg \\
Tier rate & $r^{\text{sp}}_{f,b}$ & Rate per kg in tier $b$ on flight $f$ & USD/kg \\
Tier minimum & $b_{\min}$ & $N$ tier: minimum charge floor & USD \\
\bottomrule
\end{tabular}
\end{center}

\paragraph{Spot rate cost per shipment.}
For shipment $k$ tendered on flight $f$ at the spot rate, the cost is the minimum across
applicable tiers:
\begin{equation}
  \text{cost}^{\text{spot}}(k, f) \defeq \min_{b \in B}\, r^{\text{sp}}_{f,b} \cdot \max(cw_k,\, b)
  \label{eq:spot-cost}
\end{equation}

Since $cw_k$ is a parameter (known at planning time), $\text{cost}^{\text{spot}}(k, f)$ is
pre-computable. The IATA rule that ``charging at the next tier with a lower rate may yield a
lower total cost'' is captured automatically by the min.

\subsection{Surcharge Parameters}

Per shipment per flight, the surcharge stack adds:
\begin{center}
\begin{tabular}{llp{6.5cm}l}
\toprule
Parameter & Symbol & Description & Unit \\
\midrule
Fuel surcharge & $\text{FSC}_f$ & USD/kg, monthly refresh & USD/kg \\
Security surcharge & $\text{SSC}_f$ & USD/kg & USD/kg \\
Origin THC & $\text{THC}_f^O$ & Per shipment, origin terminal handling & USD \\
Destination THC & $\text{THC}_f^D$ & Per shipment, destination terminal handling & USD \\
AMS fee (US imports) & $\text{AMS}_f$ & Per shipment & USD \\
\bottomrule
\end{tabular}
\end{center}

\subsection{Hub Connection Parameters (Through-ULD vs.\ Re-ULDing)}
\label{sec:through-uld}

For each pair of consecutive air arcs $(f, g)$ at hub $j$ — meaning $f$ ends at $j$ and $g$
starts at $j$ — and each ULD type $u$:

\begin{center}
\begin{tabular}{llp{7cm}l}
\toprule
Parameter & Symbol & Description & Unit \\
\midrule
Through-ULD compatible & $\psi_{f,g,u}$ & 1 if a ULD of type $u$ can transit hub $j$ without being broken down (same aircraft for through-flights; or same airline through-cargo agreement) & boolean \\
MCT (through) & $\text{MCT}_{f,g}^{\text{thru}}$ & Min connect time when through-ULD is used (technical-stop turnaround or hub through-cargo time) & hours \\
MCT (re-ULD) & $\text{MCT}_{f,g}^{\text{reULD}}$ & Min connect time when ULD must be broken down and rebuilt at hub & hours \\
Rehandling cost & $\rho_{f,g,u}^{\text{reULD}}$ & Cost per ULD to break down at hub and rebuild for onward flight & USD/ULD \\
\bottomrule
\end{tabular}
\end{center}

\begin{defn}[Re-ULDing]
\emph{Re-ULDing} is the operation of breaking down a built ULD at an intermediate (hub)
airport and stuffing its cargo into a different ULD for the onward flight. It occurs when:
\begin{enumerate}[noitemsep,label=(\alph*)]
  \item the two consecutive flights are operated by \emph{different airlines} (interline
        routing) — each airline owns a distinct ULD pool and ULDs cannot transfer;
  \item the two flights require \emph{different ULD types} (e.g., LD3 on a passenger belly
        leg and PMC main-deck pallet on a freighter leg);
  \item the hub lacks through-cargo facilities for this ULD type (rare at major hubs).
\end{enumerate}
Re-ULDing is operationally expensive: it adds $\sim 4$--$8$ hours of dwell time
(breakdown + storage + rebuild), incurs $\sim$\$150--\$500 in handling fees per ULD, and
introduces cargo damage risk. When $\psi_{f,g,u} = 1$, the ULD physically stays intact
across the hub — either remaining on the same aircraft (through-flight technical stop) or
moving as a complete unit between aircraft at hub through-cargo facilities.
\end{defn}

\paragraph{MVP simplification.}
In MVP, $\psi_{f,g,u}$ is a \emph{pre-computed parameter}, not a decision variable. For each
consecutive arc pair at a hub:
\begin{itemize}[noitemsep]
  \item If $\text{carrier}(f) = \text{carrier}(g)$ \emph{and} $\text{fno}(f) = \text{fno}(g)$
        (through-flight, same aircraft, intermediate stop): $\psi_{f,g,u} = 1$ for all $u \in U_f \cap U_g$.
  \item If $\text{carrier}(f) = \text{carrier}(g)$ \emph{and} the hub has a through-cargo
        agreement for type $u$: $\psi_{f,g,u} = 1$.
  \item Otherwise (interline or incompatible ULD type): $\psi_{f,g,u} = 0$, re-ULDing required.
\end{itemize}
The MILP uses $\psi$ to select the applicable MCT and rehandling cost on each used arc pair.
The choice between ``use through-ULD'' and ``re-ULD anyway'' is not a free decision in MVP;
through-ULD is used whenever $\psi = 1$. P1 extension: promote $\psi$ to a decision variable
to allow forwarder-elected re-ULDing (e.g., for re-sorting cargo destined to different
onward routes at the hub).

%%--------------------------------------------------------------------
\section{Commodity-Specific Subgraph Construction}
\label{sec:prefilter}
%%--------------------------------------------------------------------

Before the MILP is built, we construct $G(N_k, A_k)$ for each commodity $k$. The subgraph
contains only arcs that are time-feasible, cargo-type-compatible, and reachable on a complete
$o(k) \to d(k)$ path.

\paragraph{Construction algorithm.}
\begin{enumerate}
  \item \textbf{Pickup pass.} For each origin CFS $j \in N_{\text{CFS,O}}$, compute earliest
        CFS arrival:
        \[
          \tau_k(j) \defeq t_k^{\text{rdy}} + \mu_{o(k),j}^{\text{PU}}
        \]
        Arc $(o(k), j) \in A_{\text{PU}}$ is candidate iff there exists a downstream feasible path.

  \item \textbf{Origin CFS pass.} For each CFS-to-airport arc $(j, a) \in A_{\text{OC}}$ where
        $j \in N_{\text{CFS,O}}$ and $a \in N_{\text{AP,O}}$, compute earliest airport arrival:
        \[
          \tau_k(a) \defeq \tau_k(j) + \mu_{j,a}^{\text{OC}} + \beta^{\text{O}}
        \]

  \item \textbf{Flight reachability pass.} For each flight $f \in F$ (arc $(a, b) \in A_{\text{AIR}}$):
        $f$ is candidate for $k$ iff
        \begin{align*}
          &\tau_k(a) \leq \text{CGC}_f \quad &\text{(can make the cargo cutoff)}\\
          &\text{ETD}_f + \mu_f^{\text{AIR}} \leq T_k^{\text{dead}} - \text{(downstream legs)} \quad &\text{(deadline reachable)}\\
          &\tau_k = \text{General} \text{ or } \text{dgr}(f) = 1 \text{ if } \tau_k = \text{DGR}  \quad &\text{(cargo-type compatibility)}\\
          &\exists\, u \in U_f : D_k \text{ fits } u  \quad &\text{(physical fit in some ULD)}
        \end{align*}

  \item \textbf{Hub transitions pass.} For each pair of consecutive air arcs $(f, g)$ at hub $h$:
        compute $\tau_k(g_{\text{start}}) = \text{ETA}_f + \text{MCT}_{f,g}^{*}$
        where $\text{MCT}^*$ uses $\text{MCT}^{\text{thru}}$ if $\psi_{f,g,u} = 1$ for some
        $u \in U_f \cap U_g$, else $\text{MCT}^{\text{reULD}}$.

  \item \textbf{Destination CFS \& delivery pass.} For each destination airport $b \in N_{\text{AP,D}}$
        reachable: include arc $(b, c) \in A_{\text{DC}}$ where $c \in N_{\text{CFS,D}}$, then
        $(c, d(k)) \in A_{\text{FD}}$.

  \item \textbf{Reachability sweep.} Forward BFS from $o(k)$ and backward BFS from $d(k)$.
        Retain only arcs on paths reachable from both ends.
\end{enumerate}

%%--------------------------------------------------------------------
\section{Decision Variables}
%%--------------------------------------------------------------------

\begin{center}
\begin{tabular}{llp{7cm}l}
\toprule
Variable & Type & Description & Domain \\
\midrule
$x_{ij}^k$ & Binary & 1 if shipment $k$ routed on arc $(i,j) \in A^k$ & $\{0, 1\}$ \\
$y_{f,u,k}^c$ & Binary & 1 if shipment $k$ uses one ULD of type $u$ on flight $f$ under contract $c$ & $\{0, 1\}$ \\
$z_{f,u}^c$ & Integer & Number of ULDs of type $u$ claimed on flight $f$ under contract $c$ & $\mathbb{Z}_{\geq 0}$ \\
$b_{f,k}$ & Binary & 1 if shipment $k$ tendered to flight $f$ at spot rate (not via contract) & $\{0, 1\}$ \\
$C_{f,u,c}$ & Continuous & Effective chargeable weight billed for contract $c$, flight $f$, ULD type $u$ (max of actual and pivot) & $\mathbb{R}_{\geq 0}$ \\
$t_k(i)$ & Continuous & Arrival time of shipment $k$ at node $i \in N_k$ & $\mathbb{R}_{\geq 0}$ \\
\bottomrule
\end{tabular}
\end{center}

%%--------------------------------------------------------------------
\section{Constraints}
%%--------------------------------------------------------------------

\subsection*{P.1 — Flow Conservation}

For each commodity $k$ and each node $n \in N_k$:
\begin{equation}
  \sum_{j : (n,j) \in A^k} x_{nj}^k \;-\; \sum_{i : (i,n) \in A^k} x_{in}^k
  \;=\;
  \begin{cases}
    +1 & \text{if } n = o(k) \\
    -1 & \text{if } n = d(k) \\
    0  & \text{otherwise}
  \end{cases}
  \label{eq:P1}
\end{equation}

\subsection*{P.2 — ULD Volume Capacity (per flight, per ULD type, per contract)}

For each flight $f \in F$, each ULD type $u \in U_f$, each contract $c \in C_f$:
\begin{equation}
  \sum_{k \in K} v_k \cdot y_{f,u,k}^c \;\leq\; \eta \cdot V_u \cdot z_{f,u}^c
  \label{eq:P2}
\end{equation}
Total volume loaded into contract-$c$ ULDs of type $u$ on flight $f$ cannot exceed the fill-rate-capped
capacity of the claimed ULDs.

\subsection*{P.3 — ULD Weight Capacity}

For each $f, u, c$ as above:
\begin{equation}
  \sum_{k \in K} cw_k \cdot y_{f,u,k}^c \;\leq\; W_u \cdot z_{f,u}^c
  \label{eq:P3}
\end{equation}

\subsection*{P.4 — Per-Flight Contract Allocation Cap}

For each flight $f \in F$, ULD type $u$, contract $c \in C_f$:
\begin{equation}
  z_{f,u}^c \;\leq\; \text{alloc}(c, f, u)
  \label{eq:P4}
\end{equation}
Forwarder cannot claim more ULDs of type $u$ on flight $f$ than the contract allocates.

\subsection*{P.5 — Aircraft Physical Position Cap (across all contracts)}

For each flight $f \in F$, ULD type $u \in U_f$:
\begin{equation}
  \sum_{c \in C_f} z_{f,u}^c \;\leq\; P_{f,u}
  \label{eq:P5}
\end{equation}
Total ULDs of type $u$ across all contracts cannot exceed the aircraft's physical positions.

\subsection*{P.6 — Period Contract Allocation Cap}

For each contract $c$, ULD type $u$, period $t$:
\begin{equation}
  \sum_{f \in F_c \,\cap\, t} z_{f,u}^c \;\leq\; \text{rem}(c, u, t)
  \label{eq:P6}
\end{equation}
Cumulative ULDs claimed in period $t$ under contract $c$ cannot exceed the period's remaining
allocation (after prior bookings).

\subsection*{P.7 — Period Minimum Utilization (hard BSA take-or-pay)}

For each contract $c$ with $\text{hard}_c = 1$, ULD type $u$, period $t$:
\begin{equation}
  \sum_{f \in F_c \,\cap\, t} z_{f,u}^c \;\geq\; \underline{M}_{c,u,t}
  \label{eq:P7}
\end{equation}
$\underline{M}_{c,u,t}$ is the per-run minimum allocation passed in from the upstream
period-commitment model. For soft BSAs, $\underline{M}_{c,u,t} = 0$.

\subsection*{P.8 — Arc-to-ULD Linkage}

For each shipment $k$, flight $f$ with $(i,j) = f^{-1}$ in $A^k_{\text{AIR}}$:
\begin{equation}
  x_{ij}^k \;=\; \sum_{c \in C_f} \sum_{u \in U_f} y_{f,u,k}^c \;+\; b_{f,k}
  \label{eq:P8}
\end{equation}
If shipment $k$ uses arc $(i,j)$, it is loaded via exactly one contract ULD or the spot rate.

\subsection*{P.9 — Spot vs.\ Contract Exclusivity (per shipment per flight)}

Implicit in P.8: at most one of $\{y_{f,u,k}^c\}_{c,u} \cup \{b_{f,k}\}$ equals 1 when $x_{ij}^k = 1$.
Formally:
\begin{equation}
  \sum_{c \in C_f} \sum_{u \in U_f} y_{f,u,k}^c \;+\; b_{f,k} \;\leq\; 1
  \label{eq:P9}
\end{equation}

\subsection*{P.10 — Effective Chargeable Weight (pivot weight linearization)}

For each contract $c$, flight $f$, ULD type $u$ — linearizes $\max(\text{actual}, \pi \cdot z)$:
\begin{align}
  C_{f,u,c} \;&\geq\; \sum_{k \in K} cw_k \cdot y_{f,u,k}^c &&\text{(at least the actual loaded weight)} \label{eq:P10a} \\
  C_{f,u,c} \;&\geq\; \pi_{c,u} \cdot z_{f,u}^c &&\text{(at least the pivot floor)} \label{eq:P10b}
\end{align}
Since $C_{f,u,c}$ is minimized in the objective (it multiplies the rate $r_c$), at optimum
$C_{f,u,c} = \max(\sum_k cw_k \cdot y, \pi_{c,u} \cdot z)$ as required. See Section~\ref{sec:linearization} for the linearization argument.

\subsection*{P.11 — Cargo Cutoff}

For each shipment $k$ and each air arc $(i,j) \in A^k_{\text{AIR}}$ with $f = f(i,j)$:
\begin{equation}
  t_k(i) \;\leq\; \text{CGC}_f + M(1 - x_{ij}^k)
  \label{eq:P11}
\end{equation}
where $M$ is a big-M constant (use a tight bound, e.g., max planning horizon).

\subsection*{P.12 — Arrival Time Propagation (ground arcs)}

For each shipment $k$ and each ground arc $(i,j) \in A^k_{\text{PU}} \cup A^k_{\text{OC}} \cup A^k_{\text{DC}} \cup A^k_{\text{FD}}$:
\begin{equation}
  t_k(j) \;\geq\; t_k(i) + \tilde{\mu}_{ij}^{\text{G}} - M(1 - x_{ij}^k)
  \label{eq:P12}
\end{equation}
where $\tilde{\mu}_{ij}^{\text{G}} = \mu_{ij}^{\text{G}} + \beta^{\text{O}}$ for $A_{\text{OC}}$,
$\mu_{ij}^{\text{G}} + \beta^{\text{D}}$ for $A_{\text{DC}}$, else $\mu_{ij}^{\text{G}}$.

\subsection*{P.13 — Arrival Time Propagation (air arcs)}

For each shipment $k$ and each air arc $(i,j) \in A^k_{\text{AIR}}$ with $f = f(i,j)$:
\begin{equation}
  t_k(j) \;\geq\; \text{ETA}_f - M(1 - x_{ij}^k)
  \label{eq:P13}
\end{equation}

\subsection*{P.14 — Hub Minimum Connect Time}

For each shipment $k$ and each pair of consecutive air arcs $(f, g)$ at hub $h$ with both
$x_{(\cdot,h)}^k = 1$ and $x_{(h,\cdot)}^k = 1$:
\begin{equation}
  t_k(g_{\text{start}}) \;\geq\; \text{ETA}_f \;+\; \text{MCT}_{f,g}^*
  \;-\; M(2 - x_{f}^k - x_{g}^k)
  \label{eq:P14}
\end{equation}
where $\text{MCT}_{f,g}^* = \text{MCT}_{f,g}^{\text{thru}}$ if $\psi_{f,g,u} = 1$ for the
ULD type $u$ used by $k$ on both arcs, otherwise $\text{MCT}_{f,g}^{\text{reULD}}$.

\subsection*{P.15 — Deadline}

For each shipment $k$:
\begin{equation}
  t_k(d(k)) \;\leq\; T_k^{\text{dead}}
  \label{eq:P15}
\end{equation}

\subsection*{P.16 — Cargo Type Compatibility}

For each shipment $k$ with $\tau_k = \text{DGR}$ and each air arc $(i,j) \in A^k_{\text{AIR}}$:
\begin{equation}
  x_{ij}^k \;\leq\; \text{dgr}(f(i,j))
  \label{eq:P16}
\end{equation}
Analogous constraints for perishable, valuable, heavy cargo.

\subsection*{P.17 — ULD Physical Fit}

For each shipment $k$ and each $(f, u, c)$ combination where $D_k$ does not fit ULD type $u$:
\begin{equation}
  y_{f,u,k}^c \;=\; 0
  \label{eq:P17}
\end{equation}
This is a pre-processing exclusion, not a runtime constraint — fit is determined from
$D_k$ vs.\ ULD contour at subgraph construction. Bin-packing of multiple shipments inside
a single ULD is approximated by the volume constraint P.2; full 3D bin-packing is deferred.

\subsection*{P.18 — Budget Cap (optional)}

For each shipment $k$ with $B_k < \infty$:
\begin{equation}
  \text{cost}_k \;\leq\; B_k
  \label{eq:P18}
\end{equation}
where $\text{cost}_k$ is the per-shipment cost contribution from the objective.

\subsection*{P.19 — Domain}

\begin{equation}
  x_{ij}^k, y_{f,u,k}^c, b_{f,k} \in \{0,1\}, \quad
  z_{f,u}^c \in \mathbb{Z}_{\geq 0}, \quad
  C_{f,u,c}, t_k(i) \in \mathbb{R}_{\geq 0}
  \label{eq:P19}
\end{equation}

%%--------------------------------------------------------------------
\section{Objective Function}
%%--------------------------------------------------------------------

Minimize total cost across all shipments:

\begin{align}
  \min \quad
  & \sum_{k \in K} \sum_{(i,j) \in A^k_{\text{PU}} \cup A^k_{\text{OC}} \cup A^k_{\text{DC}} \cup A^k_{\text{FD}}} c_{ij}^{\text{G}} \cdot x_{ij}^k
    \tag{ground cost} \\
  +\;
  & \sum_{c \in C} \sum_{f \in F_c} \sum_{u \in U_f} r_c \cdot C_{f,u,c}
    \tag{contracted ULD cost incl.\ pivot floor} \\
  +\;
  & \sum_{k \in K} \sum_{f \in F^k} \text{cost}^{\text{spot}}(k, f) \cdot b_{f,k}
    \tag{spot rate cost} \\
  +\;
  & \sum_{k \in K} \sum_{f \in F^k} \bigl(\text{FSC}_f + \text{SSC}_f\bigr) \cdot cw_k \cdot \mathbb{1}[k \text{ on } f]
    \;+\; \text{(THC, AMS — per shipment per arc)}
    \tag{surcharges} \\
  +\;
  & \sum_{k \in K} \sum_{(f,g) \in \text{Hub}(k)} \sum_{u} \rho_{f,g,u}^{\text{reULD}} \cdot (1 - \psi_{f,g,u}) \cdot y_{f,u,k}^c \cdot y_{g,u,k}^{c'}
    \tag{rehandling cost when re-ULDing}
\end{align}

The rehandling cost term is bilinear ($y \cdot y$); linearize by introducing an auxiliary
binary variable per consecutive air-arc pair per shipment per ULD type that equals the
product (Section~\ref{sec:linearization}).

\paragraph{Indicator $\mathbb{1}[k \text{ on } f]$.}
Equals $1$ if shipment $k$ uses flight $f$; can be substituted as $\sum_{c,u} y_{f,u,k}^c + b_{f,k}$.

%%--------------------------------------------------------------------
\section{Linearization Details}
\label{sec:linearization}
%%--------------------------------------------------------------------

\subsection{Pivot Weight Linearization (P.10)}

The native cost expression per (contract, flight, ULD type) is:
\[
  \text{cost}_{f,u,c} = r_c \cdot \max\!\Bigl(\sum_k cw_k \cdot y_{f,u,k}^c,\ \pi_{c,u} \cdot z_{f,u}^c\Bigr)
\]
The $\max$ operator is non-linear. Introduce auxiliary variable $C_{f,u,c} \in \mathbb{R}_{\geq 0}$
with two linear inequality constraints (P.10a, P.10b). Replace the cost term with $r_c \cdot C_{f,u,c}$
in the objective. \textbf{Why this works:} the objective minimizes the weighted sum that includes
$r_c \cdot C_{f,u,c}$, so the optimizer pushes $C_{f,u,c}$ down to the largest binding
constraint, i.e., the maximum of the two right-hand sides. Adds one continuous variable and
two inequality constraints per (contract, flight, ULD type) triple.

\subsection{Bilinear Rehandling Cost Linearization}

The rehandling cost term contains $y_{f,u,k}^c \cdot y_{g,u,k}^{c'}$ (product of two binaries).
Standard McCormick linearization: introduce auxiliary binary $\zeta_{f,g,u,k} \in \{0,1\}$
representing the product, with three constraints:
\begin{align}
  \zeta_{f,g,u,k} \;&\leq\; y_{f,u,k}^c \\
  \zeta_{f,g,u,k} \;&\leq\; y_{g,u,k}^{c'} \\
  \zeta_{f,g,u,k} \;&\geq\; y_{f,u,k}^c + y_{g,u,k}^{c'} - 1
\end{align}
Since $\zeta$ appears with positive coefficient in the objective (rehandling cost), the
optimizer drives $\zeta$ down to 0 unless both $y$ values force it to 1 — exactly the
desired AND semantics. Adds one binary variable and three inequalities per consecutive
arc pair per ULD type per shipment.

\subsection{Big-M Tightening}

The big-M constants in P.11, P.12, P.13, P.14 can be replaced with the tightest valid
bound: e.g., the maximum possible time span across the planning horizon. Loose big-M values
weaken the LP relaxation and slow solving. Practical tightening: use $M = T_k^{\text{dead}} - t_k^{\text{rdy}} + \epsilon$
on time-related constraints.

\subsection{Spot Rate Per-Shipment Pre-Computation}

The IATA weight-break rule (\S\ref{sec:spot-params}) is a $\min$ over tiers of piecewise
linear cost. Since shipment chargeable weight $cw_k$ is a parameter, $\text{cost}^{\text{spot}}(k, f)$
is fully pre-computed before the MILP is built. No linearization needed at solve time.

%%--------------------------------------------------------------------
\section{Re-ULDing — Operational Mechanics}
\label{sec:reulding}
%%--------------------------------------------------------------------

This section motivates and defines re-ULDing in operational detail, which is the basis for
the through-ULD compatibility parameter $\psi$ and the rehandling cost $\rho^{\text{reULD}}$.

\paragraph{What re-ULDing is.}
At an intermediate (hub) airport on a multi-hop route, a built ULD arriving on flight $f$
may need to be \emph{broken down} — i.e., physically dismantled, with cargo removed
piece-by-piece — and \emph{rebuilt} into a different ULD for the onward flight $g$. The
specific sequence:
\begin{enumerate}[noitemsep]
  \item ULD arrives on inbound flight $f$ at the hub.
  \item Ground handler unloads ULD from aircraft into the hub cargo terminal.
  \item ULD is moved to the breakdown area.
  \item Cargo is removed from the ULD; each AWB's pieces are extracted and tagged.
  \item Cargo is held in the hub warehouse pending the onward flight $g$.
  \item For onward flight $g$: an empty ULD from carrier($g$)'s pool is delivered to build-up.
  \item Cargo is restuffed into the new ULD, secured with netting, weighed.
  \item Built ULD is moved to the aircraft loading area and loaded onto $g$.
\end{enumerate}

\paragraph{When re-ULDing is required ($\psi = 0$).}
\begin{itemize}[noitemsep]
  \item \textbf{Interline (different carriers):} carrier($f$) $\neq$ carrier($g$). Each
        airline maintains its own ULD pool; ULDs cannot transfer between pools. Cargo must
        leave airline $f$'s ULD and enter airline $g$'s ULD.
  \item \textbf{ULD type change:} if the inbound aircraft uses LD3s (wide-body belly) and
        the outbound aircraft uses PMC main-deck pallets (freighter), the ULD type itself
        must change. Cargo is re-stuffed into the new ULD type.
  \item \textbf{Hub lacks through-cargo facility for type $u$:} rare at major hubs (HKG,
        FRA, NRT all have through-cargo handling) but possible at smaller airports.
  \item \textbf{Forwarder elects re-ULD:} e.g., to re-sort cargo destined to different
        onward flights into different ULDs. Not modeled in MVP (deferred).
\end{itemize}

\paragraph{When through-ULD applies ($\psi = 1$).}
\begin{itemize}[noitemsep]
  \item \textbf{Through-flight technical stop:} same flight number, same aircraft. The ULD
        stays in its aircraft position throughout the intermediate stop. The hub processing
        is limited to refueling and crew change. Example: 5Y9123 TPE$\to$ANC$\to$JFK, ULD
        loaded at TPE stays on the aircraft through ANC and continues to JFK.
  \item \textbf{Same-airline through-cargo at hub:} carrier($f$) = carrier($g$), the airline
        has a through-cargo agreement at the hub, and the ULD type is compatible with both
        aircraft. ULD is unloaded from $f$, moved as a complete unit through the hub
        through-cargo facility, and loaded onto $g$. Common at major hubs for same-airline
        connections (e.g., Cathay TPE$\to$HKG$\to$JFK on CX564 + CX880).
\end{itemize}

\paragraph{Cost and time impact.}
\begin{itemize}[noitemsep]
  \item Through-ULD MCT: $\text{MCT}^{\text{thru}}$ typically $1.5$--$4$ hours depending on
        hub and aircraft turnaround.
  \item Re-ULD MCT: $\text{MCT}^{\text{reULD}}$ typically $6$--$12$ hours including
        breakdown, hub storage, and rebuild for onward flight.
  \item Rehandling cost: $\rho^{\text{reULD}}$ typically \$150--\$500 per ULD, depending on
        cargo characteristics and hub labor rates.
  \item Cargo risk: re-ULDing introduces handling-induced damage and piece-loss risk;
        valuable / fragile cargo prefers through-ULD where possible (modeled implicitly via
        the cost differential).
\end{itemize}

%%--------------------------------------------------------------------
\section{Open Items and Future Extensions}
\label{sec:deferred}
%%--------------------------------------------------------------------

\begin{enumerate}
  \item \textbf{Period-level commitment optimization (P.7).} MVP takes $\underline{M}_{c,u,t}$
        from a separate upstream model. Future: joint optimization of period commitment
        allocation across the rolling horizon.
  \item \textbf{Through-ULD as decision variable.} MVP makes $\psi$ a pre-computed
        parameter. P1: promote to decision variable, allowing forwarder-elected re-ULDing
        at compatible hubs for cargo re-sorting.
  \item \textbf{Probabilistic transit time objective.} MVP uses deterministic mean transit
        and a hard deadline. P1: extend to $\Pr(\text{arrival} \leq T_k^{\text{dead}}) \geq p$.
  \item \textbf{Full 3D bin-packing within ULD.} MVP uses aggregate volume constraint (P.2)
        with $\eta = 0.97$ as a fill-rate proxy. P1: 3D bin-packing model for shipments
        with awkward piece dimensions.
  \item \textbf{ULD availability at CFS.} MVP assumes empty ULDs are always delivered to CFS
        before cargo cutoff. P1: model ULD pool availability as a constraint.
  \item \textbf{CFS dock door capacity.} MVP assumes infinite CFS throughput. P1: model dock
        door and labor capacity as a CFS-level constraint.
  \item \textbf{Cargo type segregation within ULD.} MVP allows any compatible commodities
        to share a ULD. P1: hazmat compatibility matrix, perishable temperature segregation.
  \item \textbf{Multiple ULD types per shipment.} MVP forces a shipment to one ULD type per
        flight. P1: allow splitting across LD3 + LD7 if economically preferred.
  \item \textbf{Carrier alliance slot sharing.} MVP treats each contract independently.
        P1: model SkyTeam / Star Alliance / oneworld slot exchanges where ULD positions on
        partner carriers can substitute.
  \item \textbf{Spot capacity availability uncertainty.} MVP treats spot rate as always
        available. Real spot capacity has day-of-departure availability risk. P1: model
        spot capacity as a stochastic constraint or with explicit booking rejection events.
\end{enumerate}

\end{document}
